\documentclass[11pt,a4paper]{article}

% ---------- Packages ----------
\usepackage[utf8]{inputenc}
\usepackage[T1]{fontenc}
\usepackage{textcomp}
\usepackage{lmodern}
\usepackage{geometry}
\usepackage{microtype}
\usepackage{booktabs}
\usepackage{tabularx}
\usepackage{longtable}
\usepackage{array}
\usepackage{makecell}
\usepackage{xcolor}
\usepackage{enumitem}
\usepackage{hyperref}
\usepackage{amsmath}
\usepackage{amssymb}
\usepackage{titlesec}

% ---------- Page setup ----------
\geometry{margin=2.5cm}

\hypersetup{
    colorlinks=true,
    linkcolor=blue!60!black,
    urlcolor=blue!60!black,
    citecolor=blue!60!black
}

\setlist[itemize]{leftmargin=1.5em}
\setlength{\parskip}{0.6em}
\setlength{\parindent}{0pt}

% ---------- Custom commands ----------
\newcommand{\done}{\textcolor{green!45!black}{\(\checkmark\) Complete}}
\newcommand{\pending}{\textcolor{orange!80!black}{Pending}}
\newcommand{\phantomcomplete}{\textcolor{green!45!black}{\(\checkmark\) Complete}}
\newcommand{\code}[1]{\texttt{#1}}
\newcommand{\psnr}{\mathrm{PSNR}}
\newcommand{\ssim}{\mathrm{SSIM}}

\titleformat{\section}{\large\bfseries}{\thesection}{0.7em}{}
\titleformat{\subsection}{\normalsize\bfseries}{\thesubsection}{0.7em}{}

% ---------- Title ----------
\title{\textbf{3D Gaussian Splatting for Volumetric Fluorescence Microscopy}\\
\large Progress Report --- June 2026}

\author{
Prepared by: Syed Muhammad Asad\\
MLCV Seminar, TU Dresden\\
Repository: \code{asadrizvi64/DeepLearning-Seminar-TeamProject}, branch: \code{main}
}

\date{Report generated: 24 June 2026}

\begin{document}

\maketitle

\section*{Executive Summary}

This project investigates whether 3D Gaussian Splatting (3DGS), a representation
developed for real-time novel-view synthesis of natural scenes, can serve as a
compact and editable model for volumetric fluorescence microscopy. All core
pipeline phases are implemented and verified on synthetic phantom data. This week,
the pipeline was advanced to \emph{real} data: a single embryo frame and a short
time series from the Fluo-N3DL-DRO \emph{Drosophila} light-sheet dataset were fitted
on the TU Dresden Alpha HPC cluster using an A100 GPU.

The headline event of this reporting period was the discovery and correction of a
\textbf{region-of-interest (ROI) extraction bug} that had caused the earlier real-data
fit to model the empty interior of the embryo rather than its nucleus-bearing surface.
After the fix, the corrected runs now resolve individual nuclei, which is consistent
with the reference visualisation in Napari. A secondary rendering artifact (a small
number of oversized ``floater'' Gaussians) was also identified and removed. One
quantitative metric, the structural-similarity index (SSIM) for the 40{,}000-Gaussian
GPU fit, remains outstanding for a purely computational reason explained in
Section~\ref{sec:outstanding}.

\begin{table}[h!]
\centering
\small
\begin{tabularx}{\textwidth}{@{}p{0.30\textwidth}p{0.18\textwidth}X@{}}
\toprule
\textbf{Phase} & \textbf{Status} & \textbf{Key Result} \\
\midrule
P0 -- Environment \& infrastructure & \done & Virtual environment, SLURM scripts, and CI configured \\
P1 -- Phantom pipeline, E1--E3 & \done & Voxel-query \(\psnr = 31.7\) dB; E1 confirms alpha-blend bias \\
P2 -- Initialisation study, E2 & \done & Intensity-weighted initialisation converges fastest \\
P3 -- Fidelity/storage Pareto, E5 & \done & 3DGS reaches \(38.6\) dB at \(6.2\) KB on phantom data \\
P4 -- Export \& viewer, E6 & \done & \code{.ply}/\code{.splat} export and WebGL 3D+t viewer implemented \\
P5 -- Temporal, E7, phantom & \phantomcomplete & Shared-identity \(39.3\) dB vs.\ native-4D \(25.9\) dB \\
P5 -- Real-data temporal, cluster & \done & Corrected fit: per-frame \(30.1\)/\(32.0\)/\(32.2\) dB; nuclei resolved \\
GPU static fit, corrected ROI & \done & \(28.71\) dB on surface-nuclei ROI (40k Gaussians) \\
GPU static fit, SSIM metric & \pending & Blocked by CPU cost of full-volume query; see Sec.~\ref{sec:outstanding} \\
\bottomrule
\end{tabularx}
\end{table}

\section{Background and Motivation}

\textbf{What is 3D Gaussian Splatting?}
A scene is represented as a set of anisotropic 3D Gaussians, each carrying a centre
position, a \(3\times3\) covariance (parameterised here by a log-scale vector and a
rotation quaternion), and an amplitude. The continuous density at any point is the sum
of the Gaussians evaluated at that point. Because the representation is differentiable,
the Gaussians' parameters can be optimised by gradient descent to match a target signal.

\textbf{Why apply it to microscopy?}
A light-sheet microscopy volume is a dense 3D scalar field (fluorophore intensity).
Stored na\"ively it is large; a typical Fluo-N3DL-DRO frame is roughly \(96\) million
voxels. A Gaussian mixture promises three advantages: (i) \emph{compression}, since a
few thousand Gaussians can summarise a structured volume; (ii) an \emph{explicit,
editable} geometric model in which each Gaussian can plausibly correspond to a nucleus,
which is useful for tracking; and (iii) a path to \emph{continuous-time} (3D+t)
representations that interpolate between acquired frames.

\textbf{The central technical question} addressed early in the project (experiment E1)
is whether the rendering formula matters. Standard 3DGS composites Gaussians by
alpha-blending along camera rays, which is correct for opaque surfaces but, as shown
below, is biased for a transparent density field. The project therefore supervises the
fit by directly querying the density at voxel centres rather than by alpha-compositing.

\section{Completed Experiments (Synthetic Phantom)}

The phantom is a controlled volume of Gaussian ``blobs'' with known positions, scales,
and (for the 4D phantom) known motion. It lets us validate correctness and measure
behaviour against ground truth before touching noisy real data.

\subsection{E1 --- Integration Bias Analysis}

\textbf{Objective.}
Determine whether alpha-compositing, the standard real-time 3DGS rendering formulation,
introduces reconstruction bias when applied to a volumetric density field.

\textbf{Why a bias is expected.}
Alpha-compositing assumes each Gaussian \emph{occludes} those behind it: contributions
are attenuated by an accumulated transmittance term. For a genuinely transparent
emission field the physically correct operation is a plain \emph{sum} (line integral)
of densities, with no occlusion. When Gaussians overlap, the occlusion term
systematically removes signal that should have been added, so the error grows with
overlap.

\textbf{Finding.}
Voxel-query supervision, which evaluates the unbiased ground-truth density, achieves
\(31.7\) dB on the 15-blob phantom. The same fitted model, when scored through
alpha-compositing, collapses to \(-7.1\) dB --- a \(38.8\) dB gap. The bias increases
monotonically with Gaussian overlap. This is the empirical justification for using a
voxel-query loss throughout the project rather than the alpha-blend loss inherited from
scene reconstruction.

\begin{table}[h!]
\centering
\begin{tabular}{@{}lc@{}}
\toprule
\textbf{Supervision / scoring mode} & \textbf{Final PSNR} \\
\midrule
Voxel-query (unbiased)        & \(31.7\) dB \\
Alpha-blend projection        & \(-7.1\) dB (diverged) \\
\bottomrule
\end{tabular}
\end{table}

\subsection{E2 --- Initialisation Strategies}

Three ways of placing the initial Gaussians were benchmarked to a \(25\) dB target on
the 15-blob phantom (500 Gaussians, 1000 iterations):

\begin{itemize}
    \item \textbf{Random uniform:} Gaussians scattered uniformly. Slowest convergence
    and highest variance across random seeds, because many Gaussians start far from any
    signal and waste iterations migrating inward.
    \item \textbf{Intensity-weighted sampling:} initial centres are drawn with
    probability proportional to local intensity, so Gaussians begin near real structure.
    Fastest and most stable; \textbf{selected as the default}.
    \item \textbf{Local-maxima seeding:} one Gaussian per detected intensity peak. Fast
    when the peak count matches the true blob count, but brittle when it underestimates
    the number of nuclei (peaks merge and structure is missed).
\end{itemize}

\subsection{E5 --- Fidelity / Storage Pareto Curve}

Four representations were swept across storage budgets on the 10-blob phantom, to place
3DGS on a fidelity-versus-size trade-off curve against the natural alternatives. The
3DGS budgets were corrected this period to blob-count-appropriate sizes (5--100
Gaussians); earlier runs used 200--2000 Gaussians, which over-crowded the 10-blob
phantom and (counter-intuitively) \emph{lowered} quality, since surplus Gaussians
compete to explain the same blob and add high-frequency noise.

\begin{table}[h!]
\centering
\scriptsize
\renewcommand{\arraystretch}{1.25}
\begin{tabularx}{\textwidth}{@{}p{0.22\textwidth}XXXXX@{}}
\toprule
\textbf{Method} & \textbf{5 units} & \textbf{10 units} & \textbf{20 units} & \textbf{50 units} & \textbf{100 units} \\
\midrule
Compressed voxels &
\makecell{\(26.1\) dB\\\(0.9\) KB} & \makecell{\(32.3\) dB\\\(1.6\) KB} &
\makecell{\(38.1\) dB\\\(2.4\) KB} & \makecell{\(48.6\) dB\\\(5.3\) KB} &
\makecell{\(54.8\) dB\\\(9.4\) KB} \\
Sparse point cloud &
\makecell{\(24.0\) dB\\\(4.1\) KB} & \makecell{\(29.9\) dB\\\(7.9\) KB} &
\makecell{\(36.3\) dB\\\(11.7\) KB} & \makecell{\(43.6\) dB\\\(15.7\) KB} &
\makecell{\(51.4\) dB\\\(20.5\) KB} \\
Implicit field (MLP) &
\makecell{\(28.5\) dB\\\(1.2\) KB} & \makecell{\(41.7\) dB\\\(4.4\) KB} &
\makecell{\(45.1\) dB\\\(9.7\) KB} & \makecell{\(48.0\) dB\\\(16.9\) KB} & -- \\
3DGS (this work) &
\makecell{\(18.4\) dB\\\(2.0\) KB} & \makecell{\(21.2\) dB\\\(2.2\) KB} &
\makecell{\(24.9\) dB\\\(2.7\) KB} & \makecell{\(37.0\) dB\\\(4.0\) KB} &
\makecell{\(38.6\) dB\\\(6.2\) KB} \\
\bottomrule
\end{tabularx}
\end{table}

\textbf{Observation.}
On this \emph{sparse} synthetic phantom, compressed voxels dominate because almost the
entire volume is empty and compresses trivially; 3DGS carries a fixed per-Gaussian
overhead (position, scale, rotation, amplitude $=$ 11 floats) that is not amortised when
there are few blobs. The phantom is therefore a near-worst case for 3DGS. On real,
\emph{dense} data, where structure fills the volume and anisotropic Gaussians can each
cover an elongated nucleus, 3DGS is expected to be far more competitive --- which is
precisely what the real-data experiments in Section~\ref{sec:realdata} begin to probe.

\subsection{E6 --- Export and WebGL Viewer}

A complete export and visualisation pipeline was implemented:

\begin{itemize}
    \item \textbf{INRIA-3DGS PLY} (17 float properties per Gaussian) for compatibility
    with external tools such as SuperSplat and \code{gsplat.js}.
    \item \textbf{\code{antimatter15} \code{.splat}} binary (32 bytes per Gaussian:
    position, scale, RGBA, rotation) for the in-house viewer.
    \item \textbf{Round-trip fidelity} verified: geometry survives export and re-import
    exactly within float32 precision.
\end{itemize}

The viewer (\code{viewer/index.html}) accepts drag-and-drop of multiple \code{t*.splat}
frames and provides a slider for 3D+t scrubbing. It renders each Gaussian with a custom
\code{ShaderMaterial}: the vertex shader sizes each point from its own scale attribute,
and the fragment shader applies a radial Gaussian falloff,
\(\exp(-d^2 \times 10)\), so each splat appears as a soft blob rather than a hard disc.

\subsection{E7 --- Temporal Representation Comparison (Phantom)}

Three ways of representing a moving volume over time were benchmarked on the 4D phantom
(12 blobs, 8 frames, smooth motion, and a cell-division event at frame~4):

\begin{table}[h!]
\centering
\small
\begin{tabularx}{\textwidth}{@{}p{0.31\textwidth}cccc@{}}
\toprule
\textbf{Variant} & \textbf{Mean PSNR} & \textbf{Params} & \textbf{Storage} & \textbf{Interpolation} \\
\midrule
Shared-identity (per-frame fit) & \(39.3\) dB & 13{,}200 & \(O(T \times N)\) & No (discrete frames) \\
Native-4D (\(t_{\mathrm{center}}/t_{\mathrm{scale}}\)) & \(25.9\) dB & 3{,}900 & \(O(N)\) & Yes, \(20.1\)--\(21.9\) dB \\
Deformation field (polynomial) & \(22.8\) dB & 5{,}100 & \(O(N)\) & Yes, \(18.5\)--\(20.4\) dB \\
\bottomrule
\end{tabularx}
\end{table}

\textbf{Interpretation of the three variants.}
\emph{Shared-identity} fits one Gaussian set per frame, warm-started from the previous
frame so a given Gaussian keeps its identity over time (useful for tracking). It is the
most faithful but stores \(T\) frames' worth of parameters. \emph{Native-4D} adds a time
centre and extent to each Gaussian, so one parameter set covers the whole sequence and
can be sampled at any continuous \(t\). \emph{Deformation} keeps a canonical Gaussian set
and moves it with a low-order polynomial motion model. Shared-identity leads per-frame
fidelity by \(13.4\) dB but uses \(3.4\times\) more storage and cannot interpolate
between frames; the other two trade fidelity for an \(O(N)\) size and true continuous-time
rendering. The right choice depends on the downstream task (archival fidelity versus
sub-frame tracking).

\section{Real-Data Experiments --- Fluo-N3DL-DRO}
\label{sec:realdata}

\subsection{Dataset}

\begin{itemize}
    \item Volume size: \(125 \times 603 \times 1272\) voxels (\(z \times y \times x\)),
    uint16, LZW-compressed TIFF.
    \item Voxel size: \(0.406 \times 0.406 \times 2.03~\mu\mathrm{m}\) (approximately
    \(5\times\) anisotropy along \(z\)).
    \item Biology: at this stage the \emph{Drosophila} blastoderm is a \textbf{hollow
    shell} --- roughly \(6000\) nuclei tile the \emph{surface} of a sphere, while the
    interior is essentially empty cytoplasm. This geometry is central to the bug described
    next.
    \item 50 time frames at 30-second intervals.
\end{itemize}

\subsection{The ROI-Extraction Bug and Its Fix}
\label{sec:roi-bug}

\textbf{Symptom.}
The first real-data fits produced a featureless, uniform cloud of Gaussians in the
viewer, with no resolvable nuclei --- completely unlike the same data in Napari, which
clearly shows a surface studded with nuclei.

\textbf{Root cause.}
Fitting is performed on a cropped region of interest, not the full volume. The crop was
selected by \code{center\_roi()}, which centres a fixed-size box on the geometric centre
of the high-intensity bounding box. For a \emph{hollow} blastoderm, the bounding box
spans the whole sphere, so its geometric centre lies in the \textbf{empty interior}. The
crop therefore captured mostly signal-free cytoplasm, and the Gaussians had no nuclear
structure to lock onto.

\textbf{Fix.}
The crop is now \textbf{anchored to the top of the high-intensity bounding box} (the
polar cap of the shell) along \(z\), while remaining centred in \(x\) and \(y\). This
places the ROI on the surface layer where nuclei actually reside. The change applies to
both \code{train\_ctc.py} and \code{train\_temporal\_ctc.py} (commit \code{63c068f}).

\textbf{Confirmation on real data.}
The corrected temporal run reports its ROI as \(z\!\in\![0,24]\) --- i.e.\ anchored at the
pole --- whereas the buggy run had used \(z\!\in\![56,68]\), squarely in the interior. The
fix is therefore confirmed not just in code but in the actual crop coordinates of the
re-run.

\begin{table}[h!]
\centering
\small
\begin{tabularx}{\textwidth}{@{}XX@{}}
\toprule
\textbf{Before fix} & \textbf{After fix} \\
\midrule
\(z\)-anchor at bounding-box \emph{centre} (\(z \approx 62\)): interior cytoplasm only. &
\(z\)-anchor at bounding-box \emph{start} (\(z \approx 0\)--\(12\)): top pole and surface nuclei. \\
Viewer: uniform diffuse cloud of dots. &
Viewer: discrete nuclear clusters on the blastoderm surface. \\
\bottomrule
\end{tabularx}
\end{table}

\subsection{Corrected Static Fit (40{,}000 Gaussians, A100 GPU)}

\begin{itemize}
    \item ROI: \(64 \times 256 \times 256\) voxels, pole-anchored.
    \item Gaussians: 40{,}000; iterations: 5{,}000; batch size: 4{,}096.
    \item Wall-clock: approximately 37 minutes (2{,}211~s) on one A100.
    \item \textbf{Final PSNR: \(28.71\) dB} (SLURM job 2255712).
\end{itemize}

\textbf{Why the corrected PSNR is \emph{lower} than the buggy run (and why that is good).}
The earlier, buggy fit scored \(29.43\) dB. At first glance the fix appears to have hurt
the metric, but the opposite is true. The buggy ROI was almost flat interior cytoplasm,
which is trivially easy to reconstruct and therefore yields a deceptively high PSNR. The
corrected ROI is full of high-frequency nuclear structure, which is genuinely harder to
fit. A modest \(0.7\) dB drop while moving from ``reconstructing nothing'' to
``reconstructing real nuclei'' is the \emph{honest} and desirable outcome; the number now
measures something meaningful.

\subsection{Corrected Real-Data Temporal Reconstruction}

A 3-frame shared-identity sequence was fitted on the corrected ROI (1{,}500 Gaussians per
frame, target shape \(24\times64\times64\); SLURM job 2255711). Only three frames were
used because only \code{t000}--\code{t002} are currently staged in the cluster workspace.

\begin{table}[h!]
\centering
\small
\begin{tabular}{@{}lccc@{}}
\toprule
\textbf{Metric} & \textbf{Frame 0} & \textbf{Frame 1} & \textbf{Frame 2} \\
\midrule
Density PSNR (dB)            & 30.09 & 32.02 & 32.16 \\
Frame-to-frame consistency (dB) & --- & 26.24 & 27.60 \\
\bottomrule
\end{tabular}
\end{table}

Temporal smoothness was \(0.099\) voxels (mean acceleration of Gaussian centres),
indicating the warm-started Gaussians move only slightly between frames, as expected over
a 60-second interval. Per-frame \code{t000}--\code{t002.splat} files were exported for
3D+t scrubbing. The later frames score higher because each is warm-started from the
already-converged previous frame.

\subsection{Secondary Finding: ``Floater'' Gaussians and Their Removal}
\label{sec:floaters}

\textbf{Symptom.}
After the ROI fix, the viewer correctly showed hundreds of discrete nuclei --- but also a
single dominant glowing sphere near the centre of each frame.

\textbf{Diagnosis.}
Direct inspection of the exported \code{.splat} scale data revealed the cause. The median
Gaussian has an RMS scale of \(\approx 1.8\) voxels (nucleus-sized), and the \(99\)th
percentile is \(\approx 5\) voxels, but a \textbf{handful of outliers reach \(25\)--\(28\)
voxels}. These outliers are thin, flat discs (e.g.\ \(\approx 25\times34\times1\) voxels)
sitting at the \emph{same} position in every frame. They are classic 3DGS
``\emph{floaters}'': a few large, low-frequency Gaussians that minimise the loss by
painting the smooth background level rather than representing any nucleus. Because the
exporter maps amplitude to brightness, these large bright discs visually dominate the
render even though only 2--5 of the 1{,}500 Gaussians are involved.

\textbf{Fix.}
A small utility, \code{scripts/prune\_splat.py} (commit \code{e0f6063}), reads the
\code{.splat} binary directly and removes Gaussians whose RMS scale exceeds a threshold
(8 voxels), which sits cleanly in the gap between real nuclei (\(\le 5\)) and floaters
(\(\ge 25\)). After pruning, the maximum scale drops from \(\approx 28\) to \(\approx 7.6\)
voxels, the nucleus-sized median is untouched, and \(\approx 1{,}497\) of 1{,}500
Gaussians are retained per frame. The central blob disappears and the nuclei render
cleanly. This is a standard, defensible post-processing step in the 3DGS literature; the
\emph{principled} alternative (a scale-regularisation penalty during training) is noted as
future work in Section~\ref{sec:nextsteps}.

\section{Outstanding Item: GPU-Fit SSIM}
\label{sec:outstanding}

The one metric not yet captured is the structural-similarity index (SSIM) of the
40{,}000-Gaussian static fit. SSIM complements PSNR by measuring perceived structural
agreement rather than raw error. The blocker is purely computational: \code{compare\_fit.py}
reconstructs the full \(64\times256\times256 \approx 4.2\)~million-voxel volume by querying
every voxel against all \(40{,}000\) Gaussians (\(\sim\!10^{11}\) evaluations). On the
login node's CPU this does not complete in a reasonable time; the two manual attempts were
interrupted. The remedy is to compute the metric where the fit itself runs --- on the GPU,
inside a batch job --- or on a subsampled volume. For reference, the phantom SSIM was
\(0.9798\) (confirming the metric code is correct), and a small local real-data ROI gave
\(0.7423\); the 40k-Gaussian GPU value is expected in the same range and will be filled in
once the metric is run as a GPU batch job.

\section{Codebase and Infrastructure}

\subsection{Module Structure (\code{volsplat/})}
\begin{itemize}
    \item \code{gaussians.py}: \code{GaussianSet}, quaternion-to-rotation, and voxel-query
    rendering (the unbiased forward model from E1).
    \item \code{train.py}: \code{train\_static()}, \code{evaluate\_full()}, and point
    sampling.
    \item \code{temporal.py}: 4D phantom generator and the three temporal variants
    (shared-identity, native-4D, deformation).
    \item \code{export.py}: \code{to\_ply()}, \code{to\_splat()}, \code{read\_ply()}
    (round-trip verified).
    \item \code{ctc.py}: Cell-Tracking-Challenge loader, percentile normalisation, and the
    (now-corrected) ROI utilities.
    \item \code{losses.py}: MSE, PSNR, and a 3D-SSIM wrapper over \code{skimage}.
\end{itemize}

\subsection{HPC Scripts (\code{scripts/hpc/})}
\begin{itemize}
    \item \code{run\_ctc\_fixed.sbatch}: corrected-ROI static fit followed by
    \code{compare\_fit.py} (added this period).
    \item \code{temporal\_ctc.sbatch}: real-data temporal reconstruction job.
    \item \code{run\_compare\_fit.sh}: login-node convenience script for the SSIM metric.
    \item All scripts initialise \code{LD\_PRELOAD}/\code{LD\_LIBRARY\_PATH} before
    \code{module purge}, fixing a recurring \code{libpython} load error under
    \code{set -u} on the Alpha cluster.
\end{itemize}

\subsection{Post-processing and Viewer}
\begin{itemize}
    \item \code{scripts/prune\_splat.py}: removes oversized floater Gaussians from a
    \code{.splat} (Section~\ref{sec:floaters}).
    \item \code{viewer/index.html}: self-contained Three.js (r160) viewer with per-Gaussian
    sizing, Gaussian-falloff shading, and a 3D+t frame slider.
\end{itemize}

\section{Next Steps}
\label{sec:nextsteps}
\begin{itemize}
    \item Compute the 40k-Gaussian SSIM as a GPU batch job to close the one outstanding
    metric (Section~\ref{sec:outstanding}).
    \item Add an optional \textbf{scale-regularisation} term to \code{train\_static} so
    floaters are suppressed at the source, making the displayed Gaussians and the reported
    PSNR refer to exactly the same set.
    \item Stage additional DRO frames (\code{t001.tif}--\code{t003.tif} and beyond) in the
    workspace to extend the temporal sequence and exercise the cell-division case on real
    data.
    \item Quantify how many Gaussians per nucleus the real fit uses, to position real-data
    3DGS on the E5 storage/fidelity curve against the synthetic baselines.
    \item Prepare the final seminar presentation, including a live viewer demonstration of
    the corrected, floater-pruned embryo sequence.
\end{itemize}

\end{document}
